\documentclass[../../main.tex]{subfiles}
\begin{document}
% SOURCE p=169; printed=154; object=ch05-ex01-vector-operations; number=1; region=bottom
\begin{exercise}\label{exercise:ch05-vector-operations}
若 $\mathcal X$ 是 $K$（$=\mathbb R$ 或 $\mathbb C$）上的赋范向量空间，则加法和标量乘法分别是从 $\mathcal X\times\mathcal X$ 和 $K\times\mathcal X$ 到 $\mathcal X$ 的连续映射。此外，范数是从 $\mathcal X$ 到 $[0,\infty)$ 的连续函数；事实上，$|\|x\|-\|y\||\le\|x-y\|$。
\end{exercise}
% SUPPLEMENT: id=v14-ch05-ex01-solution; parent=ch05-ex01-vector-operations; kind=solution+motivation+hypotheses; origin=independent-derivation
\begin{solution}
\textbf{补写解答（重构新增）。}三角不等式给
\begin{align*}
\|(x+y)-(x_0+y_0)\|\le\|x-x_0\|+\|y-y_0\|,
\end{align*}
所以两坐标分别充分接近时和也接近。对标量乘法，用
\begin{align*}
\|\lambda x-\lambda_0x_0\|
\le |\lambda|\|x-x_0\|+|\lambda-\lambda_0|\|x_0\|.
\end{align*}
限制 $|\lambda-\lambda_0|<1$ 后，$|\lambda|\le|\lambda_0|+1$。任给 $\varepsilon>0$，再要求 $\|x-x_0\|<\varepsilon/[2(|\lambda_0|+1)]$ 及 $|\lambda-\lambda_0|<\varepsilon/[2(\|x_0\|+1)]$，上式便小于 $\varepsilon$，证明联合连续。

由 $\|x\|\le\|x-y\|+\|y\|$ 得 $\|x\|-\|y\|\le\|x-y\|$；交换 $x,y$ 得反向，合起来即所求反三角不等式。因此范数甚至是 Lipschitz 常数为 $1$ 的函数。
\end{solution}
\begin{remark}
\textbf{条件说明（重构新增）。}标量与向量同时变化时，须先控制标量大小，不能只证明分别连续便省去联合估计。反三角不等式把范数极限的交换归结为距离控制，全程无需完备性。
\end{remark}
% END SUPPLEMENT: id=v14-ch05-ex01-solution
% SOURCE p=169; printed=154; object=ch05-ex02-operator-norm; number=2; region=bottom
\begin{exercise}\label{exercise:ch05-operator-norm}
$L(\mathcal X,\mathcal Y)$ 是向量空间，\eqref{align:ch05-operator-norm}定义的 $\|\cdot\|$ 是其上的范数。特别地，该式右边三个表达式总相等。
\end{exercise}
% SUPPLEMENT: id=v14-ch05-ex02-solution; parent=ch05-ex02-operator-norm; kind=solution+source-clarification+motivation+hypotheses; origin=independent-derivation
\begin{solution}
\textbf{补写解答（重构新增）。}若 $S,T$ 有界线性，则 $aS+bT$ 线性，且 $\|(aS+bT)x\|\le(|a|C_S+|b|C_T)\|x\|$；零映射也有界，故成向量空间。

先设 $\mathcal X\ne\{0\}$。每个非零 $x$ 可唯一写作 $\|x\|(x/\|x\|)$，其中后一个向量范数为 $1$，故前两个上确界相等。记共同值为 $M$；它非负且由有界性而有限。对所有非零 $x$，比例至多 $M$，零向量处估计也成立，故 $M$ 本身是第三式的一个界；反过来，每个该式中的界均至少为全部比例的上确界 $M$。于是三式相等，且下确界实际在 $M$ 达到。单位闭球上的上确界也相等，因为缩放到单位球面不会减小非零像的范数。

齐次性由 $\|(aT)x\|=|a|\|Tx\|$ 后取上确界得到。逐点三角估计给 $\|(S+T)x\|\le(\|S\|+\|T\|)\|x\|$，故 $\|S+T\|\le\|S\|+\|T\|$。若 $\|T\|=0$，则 $\|Tx\|\le0$ 对每个 $x$，所以 $T=0$；反之显然。范数公理全部成立。

若 $\mathcal X=\{0\}$，则算子空间只有零映射，按正文说明采用单位闭球定义及 $C\ge0$ 的下确界，所得范数为零，范数公理同样成立。原两个空集上确界须按这一约定理解，不能在实数中直接把空集上确界当作现成的零。
\end{solution}
\begin{remark}
\textbf{动机与条件（重构新增）。}算子范数记录最佳线性误差放大率。定义域非零使球面非空；统一改用闭单位球即可覆盖零空间。有界线性两个条件分别保证上确界有限与缩放关系，缺任一项就不能原样使用这三个公式。
\end{remark}
% END SUPPLEMENT: id=v14-ch05-ex02-solution
% SOURCE p=170; printed=155; object=ch05-ex03-operator-completeness; number=3; region=top
\begin{exercise}\label{exercise:ch05-operator-completeness}
补全\cref{proposition:ch05-operator-space-complete}的证明。
\end{exercise}
% SUPPLEMENT: id=v14-ch05-ex03-solution; parent=ch05-ex03-operator-completeness; kind=solution+motivation+hypotheses; origin=independent-derivation
\begin{solution}
\textbf{补写解答（重构新增）。}设 $(T_n)$ 为算子范数 Cauchy 列。对每个 $x$，$\|T_nx-T_mx\|\le\|T_n-T_m\|\|x\|$，所以由 $\mathcal Y$ 完备得 $Tx=\lim_nT_nx$。线性组合的极限满足 $T(ax+by)=aTx+bTy$，因为两边对应序列的差范数至多 $|a|\|T_nx-Tx\|+|b|\|T_ny-Ty\|$。

Cauchy 性给 $M=\sup_n\|T_n\|<\infty$；把 $\|T_nx\|\le M\|x\|$ 取极限得 $T$ 有界。给定 $\varepsilon>0$，选 $N$ 使 $n,m\ge N$ 时 $\|T_n-T_m\|<\varepsilon$。固定任意 $x$ 及 $n\ge N$，令 $m\to\infty$ 得 $\|(T_n-T)x\|\le\varepsilon\|x\|$，所以 $\|T_n-T\|\le\varepsilon$。最后 $|\|T_n\|-\|T\||\le\|T_n-T\|\to0$，补齐所需全部结论。
\end{solution}
\begin{remark}
\textbf{条件说明（重构新增）。}逐点取极限用值域完备；在单位球上取同一个界用算子范数 Cauchy 性。定义域完备性并未参与。
\end{remark}
% END SUPPLEMENT: id=v14-ch05-ex03-solution
% SOURCE p=170; printed=155; object=ch05-ex04-evaluation; number=4; region=top
\begin{exercise}\label{exercise:ch05-evaluation-continuous}
若 $\mathcal X,\mathcal Y$ 为赋范向量空间，则映射 $(T,x)\mapsto Tx$ 从 $L(\mathcal X,\mathcal Y)\times\mathcal X$ 到 $\mathcal Y$ 连续。即若 $T_n\to T$ 且 $x_n\to x$，则 $T_nx_n\to Tx$。
\end{exercise}
% SUPPLEMENT: id=v14-ch05-ex04-solution; parent=ch05-ex04-evaluation; kind=solution+motivation+hypotheses; origin=independent-derivation
\begin{solution}
\textbf{补写解答（重构新增）。}在 $\|S-T\|<1$ 时，$\|S\|\le\|T\|+1$，而
\begin{align*}
\|Sy-Tx\|\le\|S\|\|y-x\|+\|S-T\|\|x\|.
\end{align*}
任给 $\varepsilon>0$，只要 $\|y-x\|<\varepsilon/[2(\|T\|+1)]$，且 $\|S-T\|<\min(1,\varepsilon/[2(\|x\|+1)])$，右边就小于 $\varepsilon$，证明联合连续。特别地，收敛序列满足这些最终邻域要求，得到原题的序列结论。
\end{solution}
\begin{remark}
\textbf{条件说明（重构新增）。}这里的 $T_n\to T$ 是算子范数收敛，给出统一的局部算子界；仅逐个固定输入的点态收敛，不是原题给定的拓扑。
\end{remark}
% END SUPPLEMENT: id=v14-ch05-ex04-solution
% SOURCE p=170; printed=155; object=ch05-ex05-closed-subspace; number=5; region=upper
\begin{exercise}\label{exercise:ch05-subspace-closure}
若 $\mathcal X$ 为赋范向量空间，则任何子空间的闭包仍为子空间。
\end{exercise}
% SUPPLEMENT: id=v14-ch05-ex05-solution; parent=ch05-ex05-closed-subspace; kind=solution+motivation+hypotheses; origin=independent-derivation
\begin{solution}
\textbf{补写解答（重构新增）。}设子空间为 $\mathcal M$。其闭包含 $0$。若 $x,y\in\overline{\mathcal M}$，对每个 $n$ 选 $x_n,y_n\in\mathcal M$ 使 $\|x_n-x\|<1/n$、$\|y_n-y\|<1/n$。对任意标量 $a,b$，$ax_n+by_n\in\mathcal M$，且其到 $ax+by$ 的距离至多 $(|a|+|b|)/n$，趋于零。因此 $ax+by\in\overline{\mathcal M}$，证明为子空间。
\end{solution}
\begin{remark}
\textbf{动机与条件（重构新增）。}有限线性组合与极限相容，所以补入所有可逼近的点不破坏线性结构。使用序列刻画闭包依赖这里的度量拓扑；一般拓扑向量空间须改用网或直接邻域论证。
\end{remark}
% END SUPPLEMENT: id=v14-ch05-ex05-solution
% SOURCE p=170; printed=155; object=ch05-ex06-finite-dimensional-norms; number=6; region=upper
\begin{exercise}\label{exercise:ch05-finite-dimensional-norms}
设 $\mathcal X$ 为有限维向量空间，$e_1,\ldots,e_n$ 为其一组基。
\begin{definition*}[本题的坐标范数]\phantomsection\label{definition:ch05-coordinate-l1-norm}
定义 $\|\sum_{j=1}^na_je_j\|_1=\sum_{j=1}^n|a_j|$。
\end{definition*}
\begin{enumerate}
\item\label{item:ch05-finite-dim-l1-norm} $\|\cdot\|_1$ 是 $\mathcal X$ 上的范数。
\item\label{item:ch05-finite-dim-coordinate-map} 映射 $(a_1,\ldots,a_n)\mapsto\sum_{j=1}^na_je_j$ 从取通常 Euclid 拓扑的 $K^n$ 到取 $\|\cdot\|_1$ 拓扑的 $\mathcal X$ 连续。
\item\label{item:ch05-finite-dim-sphere} $\{x\in\mathcal X:\|x\|_1=1\}$ 在 $\|\cdot\|_1$ 拓扑中紧。
\item\label{item:ch05-finite-dim-equivalence} $\mathcal X$ 上所有范数都等价。（把任意范数与 $\|\cdot\|_1$ 比较。）
\end{enumerate}
\end{exercise}
% SUPPLEMENT: id=v14-ch05-ex06-solution; parent=ch05-ex06-finite-dimensional-norms; kind=solution+motivation+hypotheses; origin=independent-derivation
\begin{solution}
\textbf{补写解答（重构新增）。}若 $n=0$，空间只有零向量，唯一范数为零，球面为空而紧，其余结论直接成立。以下设 $n\ge1$。
\begin{enumerate}
\item 对应\subcref{exercise:ch05-finite-dimensional-norms}{item:ch05-finite-dim-l1-norm}：基保证坐标唯一，所以表达式良定义；绝对值之和非负、绝对齐次，且为零当且仅当全部坐标为零。由 $|a_j+b_j|\le|a_j|+|b_j|$ 求和得三角不等式。
\item 对应\subcref{exercise:ch05-finite-dimensional-norms}{item:ch05-finite-dim-coordinate-map}：记坐标映射为 $J$，则
\begin{align*}
\|J(a)-J(b)\|_1=\sum_{j=1}^n|a_j-b_j|
\le n\left(\sum_{j=1}^n|a_j-b_j|^2\right)^{1/2},
\end{align*}
故为连续映射。
\item 对应\subcref{exercise:ch05-finite-dimensional-norms}{item:ch05-finite-dim-sphere}：$K^n$ 中 $S=\{a:\sum|a_j|=1\}$ 是闭有界集。在 $K=\mathbb C$ 时识别为 $\mathbb R^{2n}$，在实情形识别为 $\mathbb R^n$，由 Heine--Borel 定理得紧。连续像 $J(S)$ 就是所求球面，故紧。
\item 对应\subcref{exercise:ch05-finite-dimensional-norms}{item:ch05-finite-dim-equivalence}：任取另一范数 $p$，令 $C=\max_jp(e_j)>0$。由三角不等式，$p(x)\le C\|x\|_1$。又 $|p(x)-p(y)|\le p(x-y)\le C\|x-y\|_1$，故 $p$ 对 $\|\cdot\|_1$ 拓扑连续。在上述非空紧球面上，它取得最小值 $c$；球面不含零且 $p$ 为范数，故 $c>0$。对 $x\ne0$ 代入 $x/\|x\|_1$，得 $p(x)\ge c\|x\|_1$，零向量处亦成立。因此 $c\|x\|_1\le p(x)\le C\|x\|_1$。任意两范数分别与此范数比较，再合并常数，便相互等价。
\end{enumerate}
\end{solution}
\begin{remark}
\textbf{动机与条件（重构新增）。}有限维性给有限坐标估计和紧单位球面，后者把逐点严格正值提升为统一正下界；无限维一般没有这个紧性。由同样的估计，若 $T:\mathcal X\to\mathcal Y$ 线性，则 $\|Tx\|\le(\max_j\|Te_j\|)\|x\|_1$，再用范数等价可知它对任意给定范数有界，故连续；目标可以是任意赋范空间。
\end{remark}
% END SUPPLEMENT: id=v14-ch05-ex06-solution
% SOURCE p=170; printed=155; object=ch05-ex07-neumann-series; number=7; region=middle
\begin{exercise}\label{exercise:ch05-neumann-series}
设 $\mathcal X$ 为 Banach 空间。
\begin{enumerate}
\item\label{item:ch05-neumann-inverse} 若 $T\in L(\mathcal X,\mathcal X)$ 且 $\|I-T\|<1$，其中 $I$ 为恒等算子，则 $T$ 可逆；事实上，级数 $\sum_{n=0}^{\infty}(I-T)^n$ 在 $L(\mathcal X,\mathcal X)$ 中收敛到 $T^{-1}$。
\item\label{item:ch05-neumann-open} 若 $T\in L(\mathcal X,\mathcal X)$ 可逆且 $\|S-T\|<\|T^{-1}\|^{-1}$，则 $S$ 可逆。因此，可逆算子集合在 $L(\mathcal X,\mathcal X)$ 中开。
\end{enumerate}
\end{exercise}
% SUPPLEMENT: id=v14-ch05-ex07-solution; parent=ch05-ex07-neumann-series; kind=solution+motivation+hypotheses; origin=independent-derivation
\begin{solution}
\textbf{补写解答（重构新增）。}
\begin{enumerate}
\item 对应\subcref{exercise:ch05-neumann-series}{item:ch05-neumann-inverse}：令 $A=I-T$、$q=\|A\|<1$。次乘性给 $\|A^n\|\le q^n$ 对 $n\ge1$ 成立，故算子级数绝对收敛。由 $L(\mathcal X,\mathcal X)$ 完备，和 $R$ 存在。对部分和 $P_N=I+A+\cdots+A^N$，有限望远镜相消得
\begin{align*}
TP_N=(I-A)P_N=I-A^{N+1}=P_NT.
\end{align*}
由 $\|TP_N-TR\|\le\|T\|\|P_N-R\|$ 和对应右乘估计，可让 $N\to\infty$，得到 $TR=RT=I$。所以 $T$ 双射且有界逆为 $R$。
\item 对应\subcref{exercise:ch05-neumann-series}{item:ch05-neumann-open}：若 $\mathcal X\ne\{0\}$，可逆 $T$ 的逆范数正。令 $A=T^{-1}(T-S)$，则 $\|A\|\le\|T^{-1}\|\|T-S\|<1$，且 $S=T(I-A)$。前一问给 $I-A$ 可逆，故 $S^{-1}=(I-A)^{-1}T^{-1}$ 有界。以 $T$ 为中心、半径 $\|T^{-1}\|^{-1}$ 的球因此全由可逆算子组成，集合开。若 $\mathcal X=\{0\}$，唯一算子也是恒等算子且可逆；可逆集等于整个单点算子空间而开。此时不使用原式中的零范数倒数。
\end{enumerate}
\end{solution}
\begin{remark}
\textbf{动机与条件（重构新增）。}这是几何级数公式在算子代数中的形式。严格小于 $1$ 保证范数级数可和，等于 $1$ 时零算子 $T=0$ 已是不可逆的反例（空间非零）。完备性让级数的和留在代数内；只证明有限部分和恒等式还不够。第二问的乘法次序须保留，因为算子一般不交换。
\end{remark}
% END SUPPLEMENT: id=v14-ch05-ex07-solution
% SOURCE p=170; printed=155; object=ch05-ex08-measure-banach; number=8; region=middle
\begin{exercise}\label{exercise:ch05-measure-banach}
设 $(X,\mathscr M)$ 为可测空间。
\begin{definition*}[复测度空间]\phantomsection\label{definition:ch05-complex-measure-space}
令 $M(X)$ 为 $(X,\mathscr M)$ 上全体复测度所成的空间。
\end{definition*}
则 $\|\mu\|=|\mu|(X)$ 是使 $M(X)$ 成为 Banach 空间的范数。（使用\cref{theorem:ch05-banach-series-criterion}。）
\end{exercise}
% SUPPLEMENT: id=v14-ch05-ex08-solution; parent=ch05-ex08-measure-banach; kind=solution+motivation+hypotheses; origin=independent-derivation
\begin{solution}
\textbf{补写解答（重构新增）。}复测度的有限线性组合仍是复测度。由全变差性质，$|a\mu|=|a||\mu|$，$|\mu+\nu|\le|\mu|+|\nu|$，故有齐次性与三角不等式。若 $\|\mu\|=0$，则对每个 $E\in\mathscr M$，$|\mu(E)|\le|\mu|(E)\le|\mu|(X)=0$，故 $\mu=0$；反向显然。

设 $\sum_n\|\mu_n\|<\infty$。定义正测度 $\lambda=\sum_n|\mu_n|$，其可数可加性由非负双重级数换序得到，且 $\lambda(X)=\sum_n\|\mu_n\|<\infty$。每个 $\mu_n\ll\lambda$，由\cref{theorem:ch03-complex-lebesgue-radon-nikodym}，存在 $g_n\in L^1(\lambda)$ 使 $\mathrm d\mu_n=g_n\,\mathrm d\lambda$；由\cref{definition:ch03-complex-total-variation}，$|\mu_n|=|g_n|\lambda$。因此 $\sum_n\int|g_n|\,\mathrm d\lambda<\infty$。

令 $g=\sum_ng_n$ 于绝对收敛处，在其余可测零集上取零，则 $g\in L^1(\lambda)$。令 $\mathrm d\mu=g\,\mathrm d\lambda$，这是复测度，且
\begin{align*}
\left\|\mu-\sum_{n=1}^N\mu_n\right\|
=\int\left|g-\sum_{n=1}^Ng_n\right|\,\mathrm d\lambda
\le\sum_{n>N}\int|g_n|\,\mathrm d\lambda
=\sum_{n>N}\|\mu_n\|\longrightarrow0.
\end{align*}
故每个绝对收敛级数都收敛，由\cref{theorem:ch05-banach-series-criterion}得完备性。
\end{solution}
\begin{remark}
\textbf{动机与条件（重构新增）。}用全部项的全变差之和作共同控制测度，把测度级数转成同一个 $L^1$ 空间中的函数级数。假设是范数和有限，因此控制测度有限，可合法使用 Radon--Nikodym 表示；原可测空间不需要预先给定有限或 $\sigma$ 有限测度。仅对每个集合有数值收敛，不会自动给全变差范数收敛。
\end{remark}
% END SUPPLEMENT: id=v14-ch05-ex08-solution
% SOURCE p=170; printed=155; object=ch05-ex09-ck-banach; number=9; region=middle-lower
\begin{exercise}\label{exercise:ch05-ck-banach}
\begin{definition*}[$C^k([0,1])$]\phantomsection\label{definition:ch05-ck-interval}
令 $C^k([0,1])$ 为 $[0,1]$ 上所有具有直到 $k$ 阶连续导数的函数所成空间，包括端点的单侧导数。
\end{definition*}
\begin{enumerate}
\item\label{item:ch05-ck-endpoints} 若 $f\in C([0,1])$，则 $f\in C^k([0,1])$ 当且仅当它在 $(0,1)$ 上 $k$ 次连续可微，且对 $j\le k$，$\lim_{x\searrow0}f^{(j)}(x)$ 和 $\lim_{x\nearrow1}f^{(j)}(x)$ 存在。（中值定理有用。）
\item\label{item:ch05-ck-complete} $\|f\|=\sum_{j=0}^k\|f^{(j)}\|_u$ 是使 $C^k([0,1])$ 成为 Banach 空间的范数。（对 $k$ 归纳。关键是：若 $\{f_n\}\subseteq C^1([0,1])$，$f_n\to f$ 一致且 $f_n'\to g$ 一致，则 $f\in C^1([0,1])$ 且 $f'=g$。简单的证明方法是说明 $f(x)-f(0)=\int_0^xg(t)\,\mathrm dt$。）
\end{enumerate}
\end{exercise}
% SUPPLEMENT: id=v14-ch05-ex09-solution; parent=ch05-ex09-ck-banach; kind=solution+motivation+hypotheses; origin=independent-derivation
\begin{solution}
\textbf{补写解答（重构新增）。}这里 $k$ 为非负整数，端点极限指有限极限；$k=0$ 的断言就是连续函数的一致范数完备性。
\begin{enumerate}
\item 对应\subcref{exercise:ch05-ck-banach}{item:ch05-ck-endpoints}：必要性由导数在闭区间连续直接得到。反向先考虑一阶：若 $f'(x)\to L$ 当 $x\searrow0$，对实值函数，中值定理给 $(f(h)-f(0))/h=f'(\xi_h)$，其中 $0<\xi_h<h$，故单侧导数等于 $L$。复值函数分别对实部、虚部应用同一论证，不要求两者有相同中间点。右端点同理。

现在将每个内部导数 $f^{(j)}$ 用假定的有限极限延到端点。$f$ 的已有连续性使零阶端点值一致。先对 $f$ 用一阶论证，再对已经延好的 $f'$、$f''$ 依次使用，便知这些延拓确为逐阶单侧导数，且各阶都连续，得到 $C^k$。
\item 对应\subcref{exercise:ch05-ck-banach}{item:ch05-ck-complete}：范数非负、齐次和三角不等式逐项成立，零阶项保证正定。设 $(f_n)$ 为此范数 Cauchy 列，则对每个 $0\le j\le k$，$(f_n^{(j)})$ 在 $C([0,1])$ 的一致范数中 Cauchy，故一致收敛到连续函数 $g_j$。对 $j<k$，微积分基本定理给
\begin{align*}
f_n^{(j)}(x)-f_n^{(j)}(0)=\int_0^xf_n^{(j+1)}(t)\,\mathrm dt.
\end{align*}
积分误差至多 $\|f_n^{(j+1)}-g_{j+1}\|_u$，故取极限得 $g_j(x)-g_j(0)=\int_0^xg_{j+1}(t)\,\mathrm dt$。因 $g_{j+1}$ 连续，基本定理说明 $g_j$ 的导数是 $g_{j+1}$，端点按单侧理解。因此 $g_0\in C^k$ 且 $g_0^{(j)}=g_j$。有限个一致误差之和趋于零，给 $\|f_n-g_0\|\to0$。
\end{enumerate}
\end{solution}
\begin{remark}
\textbf{动机与逐条件说明（重构新增）。}范数同时控制函数及全部导数，才能在极限中保留微分关系。只让函数一致收敛，极限未必可微。端点有限极限与零阶连续性共同保证内部导数接到正确端点值；若允许导数极限无穷大，$\sqrt x$ 就不能在 $0$ 有有限单侧导数。
\end{remark}
% END SUPPLEMENT: id=v14-ch05-ex09-solution
% SOURCE p=170; printed=155; object=ch05-ex10-sobolev-banach; number=10; region=lower
\begin{exercise}\label{exercise:ch05-sobolev-banach}
\begin{definition*}[$L_k^1([0,1])$]\phantomsection\label{definition:ch05-l1k-interval}
令 $L_k^1([0,1])$ 为所有 $f\in C^{k-1}([0,1])$ 且 $f^{(k-1)}$ 在 $[0,1]$ 上绝对连续的函数所成空间；因而 $f^{(k)}$ 几乎处处存在并属于 $L^1([0,1])$。
\end{definition*}
则 $\|f\|=\sum_{j=0}^k\int_0^1|f^{(j)}(x)|\,\mathrm dx$ 是使 $L_k^1([0,1])$ 成为 Banach 空间的范数。（见\cref{exercise:ch05-ck-banach}及其提示。）
\end{exercise}
% SUPPLEMENT: id=v14-ch05-ex10-solution; parent=ch05-ex10-sobolev-banach; kind=solution+motivation+hypotheses; origin=independent-derivation
\begin{solution}
\textbf{补写解答（重构新增）。}此定义取正整数 $k\ge1$。先注意对任意绝对连续函数 $u$，由积分表示，对 $x,y\in[0,1]$ 有 $|u(x)|\le|u(y)|+\|u'\|_1$。固定 $x$ 对 $y$ 积分，得到关键估计
\begin{align*}
\|u\|_u\le\|u\|_1+\|u'\|_1.
\end{align*}
复值情形使用复值绝对连续函数的积分表示，估计仍成立。

该函数类对线性组合封闭。范数的齐次性和三角不等式逐项由积分得到；若范数为零，零阶项给 $f=0$ 几乎处处，而 $f$ 连续，若某点不为零则其某邻域上绝对值有正下界，积分不可能为零，故处处为零。因此确为范数。

设 $(f_n)$ 为 Cauchy 列。对 $0\le j<k$，$f_n^{(j)}-f_m^{(j)}$ 绝对连续：最高一阶由定义，较低阶是 $C^1$ 函数。由上述估计，$(f_n^{(j)})$ 在一致范数中 Cauchy，故一致趋于连续函数 $g_j$；最高阶 $(f_n^{(k)})$ 在 $L^1$ 中 Cauchy，故趋于某 $g_k\in L^1$。

对每个 $j<k$，积分表示
\begin{align*}
f_n^{(j)}(x)-f_n^{(j)}(0)=\int_0^xf_n^{(j+1)}(t)\,\mathrm dt
\end{align*}
均成立。较低阶的被积函数一致收敛，最后一阶在 $L^1$ 中收敛，两种情形的积分误差都对 $x$ 一致趋零。因此
\begin{align*}
g_j(x)-g_j(0)=\int_0^xg_{j+1}(t)\,\mathrm dt\qquad(0\le j<k).
\end{align*}
对 $j<k-1$，被积函数连续，所以 $g_j'=g_{j+1}$；对 $j=k-1$，被积函数可积，所以 $g_{k-1}$ 绝对连续，导数几乎处处为 $g_k$。于是 $f=g_0\in L_k^1$。低阶导数的一致收敛蕴含 $L^1$ 收敛，最高阶已有 $L^1$ 收敛，有限求和给 $\|f_n-f\|\to0$，证明完备。
\end{solution}
\begin{remark}
\textbf{动机与条件（重构新增）。}绝对连续性提供精确积分表示，使最高阶导数的 $L^1$ 控制可以向下传回函数。仅几乎处处可微且导数可积并不足够，Cantor 函数就是差别的例子。区间有限且长度为 $1$，使积分平均给出上述统一估计。定义取连续代表而最高阶导数只按几乎处处等同，二者的角色没有混淆。
\end{remark}
% END SUPPLEMENT: id=v14-ch05-ex10-solution
% SOURCE p=170,171; printed=155,156; object=ch05-ex11-holder-banach; number=11; region=bottom-top
\begin{exercise}\label{exercise:ch05-holder-banach}
\begin{definition*}[Hölder 空间与小 Hölder 类]\phantomsection\label{definition:ch05-holder-spaces}
若 $0<\alpha\le1$，令 $\Lambda_\alpha([0,1])$ 为指数 $\alpha$ 的 Hölder 连续函数空间，即 $f\in\Lambda_\alpha$ 当且仅当
\begin{align*}
\|f\|_{\Lambda_\alpha}=|f(0)|+\sup_{\substack{x,y\in[0,1]\\x\ne y}}\frac{|f(x)-f(y)|}{|x-y|^\alpha}<\infty.
\end{align*}
令 $\lambda_\alpha([0,1])$ 为其中对每个 $y\in[0,1]$ 都满足 $|f(x)-f(y)|/|x-y|^\alpha\to0$（$x\to y$）的函数集合。
\end{definition*}
\begin{enumerate}
\item\label{item:ch05-holder-complete} $\|\cdot\|_{\Lambda_\alpha}$ 是使 $\Lambda_\alpha([0,1])$ 成为 Banach 空间的范数。
\item\label{item:ch05-little-holder} 若 $\alpha<1$，则 $\lambda_\alpha([0,1])$ 是 $\Lambda_\alpha([0,1])$ 的无限维闭子空间；若 $\alpha=1$，则 $\lambda_\alpha([0,1])$ 只含常函数。
\end{enumerate}
\end{exercise}
% SUPPLEMENT: id=v14-ch05-ex11-solution; parent=ch05-ex11-holder-banach; kind=solution+motivation+hypotheses; origin=independent-derivation
\begin{solution}
\textbf{补写解答（重构新增）。}记上式差商上确界为 $[f]_\alpha$。
\begin{enumerate}
\item 对应\subcref{exercise:ch05-holder-banach}{item:ch05-holder-complete}：齐次性和三角不等式分别对基点项与差商项成立。若范数为零，则 $[f]_\alpha=0$ 使 $f$ 为常数，而 $f(0)=0$ 使该常数为零。因此为范数。并且 $\|f\|_u\le|f(0)|+[f]_\alpha=\|f\|_{\Lambda_\alpha}$，因为 $|x|^\alpha\le1$。

若 $(f_n)$ 为该范数 Cauchy 列，则一致趋于某连续 $f$。给定 $\varepsilon>0$，选 $N$ 使 $n,m\ge N$ 时 $\|f_n-f_m\|_{\Lambda_\alpha}<\varepsilon$。固定 $n\ge N$ 及任意不同 $x,y$，令 $m\to\infty$ 得
\begin{align*}
|f_n(0)-f(0)|+\frac{|(f_n-f)(x)-(f_n-f)(y)|}{|x-y|^\alpha}\le\varepsilon.
\end{align*}
对点对取上确界，得 $\|f_n-f\|_{\Lambda_\alpha}\le\varepsilon$；特别 $f-f_n$ 有有限范数，故 $f\in\Lambda_\alpha$，并得到范数收敛。
\item 对应\subcref{exercise:ch05-holder-banach}{item:ch05-little-holder}：对固定 $y$ 的零极限条件由三角不等式和齐次性对线性组合保持，所以为子空间。若 $f_n\in\lambda_\alpha$ 且 $f_n\to f$ 于 $\Lambda_\alpha$，固定 $y$ 及 $\varepsilon>0$，先选 $n$ 使 $[f-f_n]_\alpha<\varepsilon/2$，再用 $f_n$ 在这个 $y$ 的零极限取邻域，使其差商小于 $\varepsilon/2$。二者相加给 $f$ 的差商小于 $\varepsilon$，所以子空间闭。这里无需先假定对 $y$ 一致的极限。

若 $\alpha<1$，每个多项式 $p$ 都满足 $|p(x)-p(y)|\le\|p'\|_u|x-y|$，故差商至多 $\|p'\|_u|x-y|^{1-\alpha}\to0$，且其 Hölder 范数有限。因此 $1,x,x^2,\ldots$ 这组线性无关函数全在其中，空间无限维。若 $\alpha=1$，零差商条件使每个内部点的导数存在且为零；实、虚部分别由中值定理在 $(0,1)$ 上恒定，连续性再延至端点。常函数显然满足条件，所以恰为常函数空间。
\end{enumerate}
\end{solution}
\begin{remark}
\textbf{动机与条件（重构新增）。}差商只看增量，无法控制常数方向，故范数加 $|f(0)|$。指数 $\alpha>0$ 使有限差商蕴含连续性；当 $\alpha<1$，可微函数的一级增量除以 $|x-y|^\alpha$ 仍趋零，而 $\alpha=1$ 恰把条件变成导数为零，解释了结论的突变。原题的小类按每个固定 $y$ 定义，证明没有擅自换成更强的一致定义。
\end{remark}
% END SUPPLEMENT: id=v14-ch05-ex11-solution
% SOURCE p=171; printed=156; object=ch05-ex12-quotient-norm; number=12; region=upper
\begin{exercise}\label{exercise:ch05-quotient-norm}
设 $\mathcal X$ 是赋范向量空间，$\mathcal M$ 是它的真闭子空间。
\begin{enumerate}
\item\label{item:ch05-quotient-is-norm} $\|x+\mathcal M\|=\inf\{\|x+y\|:y\in\mathcal M\}$ 是 $\mathcal X/\mathcal M$ 上的范数。
\item\label{item:ch05-quotient-almost-unit} 对任意 $\varepsilon>0$，存在 $x\in\mathcal X$ 使 $\|x\|=1$ 且 $\|x+\mathcal M\|\ge1-\varepsilon$。
\item\label{item:ch05-quotient-projection-norm} 投影 $\pi(x)=x+\mathcal M$ 从 $\mathcal X$ 到 $\mathcal X/\mathcal M$ 的范数为 $1$。
\item\label{item:ch05-quotient-complete} 若 $\mathcal X$ 完备，则 $\mathcal X/\mathcal M$ 也完备。（使用\cref{theorem:ch05-banach-series-criterion}。）
\item\label{item:ch05-quotient-topology} 商范数定义的拓扑就是\cref{exercise:ch04-quotient-topology}所定义的商拓扑。
\end{enumerate}
\end{exercise}
% SUPPLEMENT: id=v14-ch05-ex12-solution; parent=ch05-ex12-quotient-norm; kind=solution+motivation+hypotheses; origin=independent-derivation
\begin{solution}
\textbf{补写解答（重构新增）。}记 $q(x+\mathcal M)=\inf_{m\in\mathcal M}\|x+m\|$。
\begin{enumerate}
\item 对应\subcref{exercise:ch05-quotient-norm}{item:ch05-quotient-is-norm}：更换代表 $x$ 为 $x+m_0$ 只是重新排列 $\mathcal M$ 中的平移项，所以下确界不变。取 $m=0$ 得 $0\le q(x+\mathcal M)\le\|x\|<\infty$。对非零标量 $a$，$m\mapsto am$ 为 $\mathcal M$ 的双射，故 $q(ax+\mathcal M)=|a|q(x+\mathcal M)$；$a=0$ 时直接成立。给定 $\eta>0$，分别选 $m,n$ 使 $\|x+m\|<q(x+\mathcal M)+\eta$、$\|y+n\|<q(y+\mathcal M)+\eta$，得到
\begin{align*}
q(x+y+\mathcal M)\le\|(x+m)+(y+n)\|
<q(x+\mathcal M)+q(y+\mathcal M)+2\eta.
\end{align*}
让 $\eta\downarrow0$ 得三角不等式。若 $q(x+\mathcal M)=0$，可选 $m_n$ 使 $x+m_n\to0$，所以 $-m_n\to x$；$\mathcal M$ 闭给 $x\in\mathcal M$，即商零元。反向由取 $m=-x$ 得到。
\item 对应\subcref{exercise:ch05-quotient-norm}{item:ch05-quotient-almost-unit}：先设 $0<\varepsilon<1$。取 $z\notin\mathcal M$，由前问 $d=q(z+\mathcal M)>0$。因 $d/(1-\varepsilon)>d$，能选 $m\in\mathcal M$ 使 $\|z+m\|<d/(1-\varepsilon)$。令 $x=(z+m)/\|z+m\|$，则 $\|x\|=1$，且 $q(x+\mathcal M)=d/\|z+m\|>1-\varepsilon$。若 $\varepsilon\ge1$，任取单位向量即可；真子空间假设保证 $\mathcal X$ 非零。
\item 对应\subcref{exercise:ch05-quotient-norm}{item:ch05-quotient-projection-norm}：$\pi$ 线性，且 $\|\pi x\|\le\|x\|$，故 $\|\pi\|\le1$。前问对任意正误差给单位向量使像范数至少 $1-\varepsilon$，故 $\|\pi\|\ge1$，从而等于 $1$。
\item 对应\subcref{exercise:ch05-quotient-norm}{item:ch05-quotient-complete}：设商空间中的级数 $\sum_nu_n$ 绝对收敛。为每个类 $u_n$ 选代表 $x_n$，使 $\pi x_n=u_n$ 且 $\|x_n\|<\|u_n\|+2^{-n}$，这是下确界的定义。于是 $\sum_n\|x_n\|<\infty$，在完备的 $\mathcal X$ 中收敛到某 $x$。由 $\pi$ 连续，$\sum_{n=1}^Nu_n=\pi(\sum_{n=1}^Nx_n)\to\pi x$。商中的每个绝对收敛级数都收敛，由判别定理得完备。
\item 对应\subcref{exercise:ch05-quotient-norm}{item:ch05-quotient-topology}：对每个 $r>0$ 有
\begin{align*}
\pi(B_{\mathcal X}(x,r))=B_{\mathcal X/\mathcal M}(\pi x,r).
\end{align*}
正包含由收缩估计；反包含来自 $q(z-x+\mathcal M)<r$ 时能选 $m\in\mathcal M$ 使 $\|z-x+m\|<r$，因此同类代表 $z+m$ 在原球中。于是 $\pi$ 为开映射。若 $V$ 商范数开，则由连续性 $\pi^{-1}(V)$ 开；反过来，若 $\pi^{-1}(V)$ 开，则 $V=\pi(\pi^{-1}(V))$ 开。恰为商拓扑的定义。
\end{enumerate}
\end{solution}
\begin{remark}
\textbf{动机与逐条件说明（重构新增）。}闭性只在半范数的正定性处不可省；真子空间条件保证存在非零商方向，从而投影范数恰为 $1$。若 $\mathcal M=\mathcal X$，商仍完备且上述拓扑结论仍成立，但投影范数为零，第二、三问不能原样保留。商完备性用近乎最短代表提升可和级数，完全不要求每个类有真正最短代表。
\end{remark}
% END SUPPLEMENT: id=v14-ch05-ex12-solution
% SOURCE p=171; printed=156; object=ch05-ex13-seminorm-nullspace; number=13; region=middle
\begin{exercise}\label{exercise:ch05-seminorm-nullspace}
若 $\|\cdot\|$ 为向量空间 $\mathcal X$ 上的半范数，令 $\mathcal M=\{x\in\mathcal X:\|x\|=0\}$。则 $\mathcal M$ 为子空间，且映射 $x+\mathcal M\mapsto\|x\|$ 是 $\mathcal X/\mathcal M$ 上的范数。
\end{exercise}
% SUPPLEMENT: id=v14-ch05-ex13-solution; parent=ch05-ex13-seminorm-nullspace; kind=solution+motivation+hypotheses; origin=independent-derivation
\begin{solution}
\textbf{补写解答（重构新增）。}记半范数为 $p$。由 $p(0)=0$，零向量在 $\mathcal M$。若 $x,y\in\mathcal M$，则 $p(ax+by)\le|a|p(x)+|b|p(y)=0$，故为子空间。若 $x-y\in\mathcal M$，三角不等式给 $p(x)\le p(y)+p(x-y)=p(y)$；交换 $x,y$ 得相反不等式，故 $p(x)=p(y)$，所以商上的表达式良定义。

齐次性和三角不等式直接从代表元下降；非负性也保留。商范数为零当且仅当代表 $x\in\mathcal M$，也就是该类为商零元。因此它是范数。
\end{solution}
\begin{remark}
\textbf{动机与条件（重构新增）。}把所有大小为零的方向除去，正好修复半范数唯一缺少的正定性。此题没有预先给拓扑，因而不另加“$\mathcal M$ 闭”的假设；等价类上的值不依赖代表是本构造的关键。
\end{remark}
% END SUPPLEMENT: id=v14-ch05-ex13-solution
% SOURCE p=171; printed=156; object=ch05-ex14-nonclosed-quotient; number=14; region=middle
\begin{exercise}\label{exercise:ch05-nonclosed-quotient}
若 $\mathcal X$ 为赋范向量空间，$\mathcal M$ 为非闭子空间，则\cref{exercise:ch05-quotient-norm}定义的 $\|x+\mathcal M\|$ 是 $\mathcal X/\mathcal M$ 上的半范数。若像\cref{exercise:ch05-seminorm-nullspace}那样再除去它的零空间，所得商空间与 $\mathcal X/\overline{\mathcal M}$ 等距同构。（参见\cref{exercise:ch05-subspace-closure}。）
\end{exercise}
% SUPPLEMENT: id=v14-ch05-ex14-solution; parent=ch05-ex14-nonclosed-quotient; kind=solution+motivation+hypotheses; origin=independent-derivation
\begin{solution}
\textbf{补写解答（重构新增）。}习题中的下确界不依赖代表、非负、齐次及满足三角不等式的证明均未使用闭性，故仍为半范数 $q$。而 $q(x+\mathcal M)=0$ 当且仅当 $x\in\overline{\mathcal M}$，由距离为零和序列逼近直接得到。它的零空间因而为 $\mathcal N=\overline{\mathcal M}/\mathcal M$，其中闭包是子空间。

定义
\begin{align*}
\Phi:(\mathcal X/\mathcal M)/\mathcal N\longrightarrow\mathcal X/\overline{\mathcal M},
\qquad (x+\mathcal M)+\mathcal N\longmapsto x+\overline{\mathcal M}.
\end{align*}
两个左端代表相同，当且仅当原代表之差在 $\overline{\mathcal M}$，正好等价于右端相同。因此映射良定义且单射；每个右端类显然有原像，故满射，并由代表运算可见线性。

最后，$\inf_{m\in\mathcal M}\|x+m\|=\inf_{z\in\overline{\mathcal M}}\|x+z\|$：右边至多左边；反向，对每个 $z\in\overline{\mathcal M}$ 取 $m_n\to z$，范数连续性给左边至多 $\lim_n\|x+m_n\|=\|x+z\|$，再对 $z$ 取下确界。这个等式表明 $\Phi$ 保持范数，故为等距同构。
\end{solution}
\begin{remark}
\textbf{动机与条件（重构新增）。}非闭子空间外仍可能有距离为零的点，它们恰是新增的闭包方向。先代数商再消去零距离，等价于一开始直接除以闭包；非闭假设描述需要修复的情形，证明也涵盖本来已经闭的情形。
\end{remark}
% END SUPPLEMENT: id=v14-ch05-ex14-solution
% SOURCE p=171; printed=156; object=ch05-ex15-kernel-factorization; number=15; region=middle-lower
\begin{exercise}\label{exercise:ch05-kernel-factorization}
设 $\mathcal X,\mathcal Y$ 为赋范向量空间，$T\in L(\mathcal X,\mathcal Y)$。
\begin{definition*}[本题的零空间]\phantomsection\label{definition:ch05-operator-kernel}
令 $\mathcal N(T)=\{x\in\mathcal X:Tx=0\}$。
\end{definition*}
\begin{enumerate}
\item\label{item:ch05-kernel-closed} $\mathcal N(T)$ 是 $\mathcal X$ 的闭子空间。
\item\label{item:ch05-kernel-factor} 存在唯一的 $S\in L(\mathcal X/\mathcal N(T),\mathcal Y)$ 使 $T=S\circ\pi$，其中 $\pi$ 为商投影（见\cref{exercise:ch05-quotient-norm}）。而且 $\|S\|=\|T\|$。
\end{enumerate}
\end{exercise}
% SUPPLEMENT: id=v14-ch05-ex15-solution; parent=ch05-ex15-kernel-factorization; kind=solution+motivation+hypotheses; origin=independent-derivation
\begin{solution}
\textbf{补写解答（重构新增）。}
\begin{enumerate}
\item 对应\subcref{exercise:ch05-kernel-factorization}{item:ch05-kernel-closed}：核含零且对线性组合封闭，由 $T$ 线性得到；又 $T$ 连续，目标度量空间的 $\{0\}$ 闭，所以核为其闭原像。
\item 对应\subcref{exercise:ch05-kernel-factorization}{item:ch05-kernel-factor}：必须令 $S(x+\mathcal N(T))=Tx$。若两个代表相差一个核向量，像相等，所以良定义；它显然线性，满足所需分解。对任意 $m\in\mathcal N(T)$，
\begin{align*}
\|S(x+\mathcal N(T))\|=\|T(x+m)\|\le\|T\|\|x+m\|.
\end{align*}
对 $m$ 取下确界，得 $\|S\|\le\|T\|$。反过来，$\|Tx\|=\|S\pi x\|\le\|S\|\|\pi x\|\le\|S\|\|x\|$，故 $\|T\|\le\|S\|$。两式给范数相等。每个商类都是某个 $\pi x$，所以分解式强制确定所有 $S$ 的值，唯一性得证。若 $T=0$，核为整个空间，商为零空间，上述公式仍给唯一零算子和范数相等，无须使用“真子空间时投影范数为1”。
\end{enumerate}
\end{solution}
\begin{remark}
\textbf{动机与条件（重构新增）。}除去核，把 $T$ 完全看不见的方向先消掉，得到一个不增加最佳界的商映射。线性性保证核结构和良定义，有界性保证核闭及商算子有界。结论不声称 $S$ 的逆有界，也不声称它满射到 $\mathcal Y$。
\end{remark}
% END SUPPLEMENT: id=v14-ch05-ex15-solution
\Needspace{18\baselineskip}
% SOURCE p=171; printed=156; object=ch05-ex16-vector-integration; number=16; region=lower
\begin{exercise}\label{exercise:ch05-vector-integration}
本题的目的是建立取值于可分 Banach 空间的函数的积分理论。设 $(X,\mathscr M,\mu)$ 为测度空间，$\mathcal Y$ 为可分 Banach 空间。
\begin{definition*}[本题的向量值函数类]\phantomsection\label{definition:ch05-vector-integrable-classes}
令 $L_{\mathcal Y}$ 为所有从 $X$ 到 $\mathcal Y$ 的 $(\mathscr M,\mathcal B_{\mathcal Y})$ 可测映射所成空间；$F_{\mathcal Y}$ 为形如 $f(x)=\sum_{j=1}^n\chi_{E_j}(x)y_j$ 的映射集合，其中 $n\in\mathbb N$，$y_j\in\mathcal Y$，$E_j\in\mathscr M$，且 $\mu(E_j)<\infty$。若 $f\in L_{\mathcal Y}$，由于 $y\mapsto\|y\|$ 连续（\cref{exercise:ch05-vector-operations}），$x\mapsto\|f(x)\|$ 是 $(\mathscr M,\mathcal B_{\mathbb R})$ 可测函数，定义
\begin{align*}
\|f\|_1=\int\|f(x)\|\,\mathrm d\mu(x),\qquad
L_{\mathcal Y}^1=\{f\in L_{\mathcal Y}:\|f\|_1<\infty\}.
\end{align*}
\end{definition*}
\begin{enumerate}
\item\label{item:ch05-vector-integral-spaces} $L_{\mathcal Y}$ 为向量空间，$F_{\mathcal Y}$ 和 $L_{\mathcal Y}^1$ 为其子空间，$F_{\mathcal Y}\subseteq L_{\mathcal Y}^1$；$\|\cdot\|_1$ 是 $L_{\mathcal Y}^1$ 上的半范数，把几乎处处相等的函数等同后它成为范数。
\item\label{item:ch05-vector-relative-balls} 设 $\{y_n\}_{n=1}^{\infty}$ 为 $\mathcal Y$ 中的可数稠密集。给定 $\varepsilon>0$，令 $B_n^\varepsilon=\{y\in\mathcal Y:\|y-y_n\|<\varepsilon\|y_n\|\}$。则 $\bigcup_{n=1}^{\infty}B_n^\varepsilon\supseteq\mathcal Y\setminus\{0\}$。
\item\label{item:ch05-vector-simple-approximation} 若 $f\in L_{\mathcal Y}^1$，则存在 $\{h_n\}\subseteq F_{\mathcal Y}$ 使 $h_n\to f$ 几乎处处且 $\|h_n-f\|_1\to0$。（用上一问的记号，令 $A_{nj}=B_n^{1/j}\setminus\bigcup_{m=1}^{n-1}B_m^{1/j}$、$E_{nj}=f^{-1}(A_{nj})$，并考虑 $g_j=\sum_{n=1}^{\infty}y_n\chi_{E_{nj}}$。）
\item\label{item:ch05-vector-integral-existence} 存在唯一的线性映射 $\int:L_{\mathcal Y}^1\to\mathcal Y$，满足对 $y\in\mathcal Y$、$E\in\mathscr M$ 且 $\mu(E)<\infty$ 有 $\int y\chi_E=\mu(E)y$，并且 $\|\int f\|\le\|f\|_1$。
% SOURCE p=172; printed=157; object=ch05-ex16-vector-integration-continuation; region=top
\item\label{item:ch05-vector-dct} 控制收敛定理：若 $\{f_n\}\subseteq L_{\mathcal Y}^1$，$f_n\to f$ 几乎处处，并存在 $g\in L^1$ 使对所有 $n$ 及几乎处处的 $x$ 有 $\|f_n(x)\|\le g(x)$，则 $\int f_n\to\int f$。
\item\label{item:ch05-vector-integral-operator} 若 $\mathcal Z$ 为可分 Banach 空间，$T\in L(\mathcal Y,\mathcal Z)$，且 $f\in L_{\mathcal Y}^1$，则 $T\circ f\in L_{\mathcal Z}^1$，并且 $\int T\circ f=T(\int f)$。
\end{enumerate}
\end{exercise}
% SUPPLEMENT: id=v14-ch05-ex16-solution; parent=ch05-ex16-vector-integration; kind=solution+source-clarification+motivation+hypotheses; origin=independent-derivation
\begin{remark}
\textbf{几乎处处极限的约定（重构新增）。}在\subcref{exercise:ch05-vector-integration}{item:ch05-vector-dct}中，$f$ 应取可测代表，或理解为几乎处处相等的等价类；若原测度不完备，任意修改零集内的值未必还可测。下文会构造可测代表再定义积分。这个约定避免额外偷偷假设 $\mu$ 完备。
\end{remark}
\begin{solution}
\textbf{逐项完整补写（重构新增）。}
\begin{enumerate}
\item 对应\subcref{exercise:ch05-vector-integration}{item:ch05-vector-integral-spaces}：可分度量空间 $\mathcal Y$ 有可数基，$\mathcal Y\times\mathcal Y$ 的开集是可数个基矩形的并。若 $f,g$ 可测，则 $(f,g)$ 对每个基矩形的原像是两个可测原像之交，故为可测映射。与连续的 $(u,v)\mapsto au+bv$ 复合，得 $af+bg$ 可测。因此 $L_{\mathcal Y}$ 为向量空间。

有限值函数的各值原像可由有限个 $E_j$ 的布尔组合表示，故 $F_{\mathcal Y}$ 中的映射可测；合并有限表示即可见它对线性组合封闭。由点态三角不等式，
\begin{align*}
\left\|\sum_{j=1}^ny_j\chi_{E_j}\right\|_1
\le\sum_{j=1}^n\|y_j\|\mu(E_j)<\infty,
\end{align*}
所以 $F_{\mathcal Y}\subseteq L_{\mathcal Y}^1$。积分三角不等式和齐次性给 $L_{\mathcal Y}^1$ 为子空间并给出半范数公理。非负函数积分为零当且仅当它几乎处处为零，所以 $\|f\|_1=0$ 当且仅当 $f=0$ 几乎处处。等价类上的运算和半范数不依赖代表；除去该零空间后范数正定。
\item 对应\subcref{exercise:ch05-vector-integration}{item:ch05-vector-relative-balls}：固定 $y\ne0$。选 $0<\delta<\varepsilon/(1+\varepsilon)$，再用稠密性取 $y_n$ 使 $\|y-y_n\|<\delta\|y\|$。则 $\|y_n\|>(1-\delta)\|y\|>0$，且
\begin{align*}
\|y-y_n\|<\delta\|y\|<\varepsilon(1-\delta)\|y\|<\varepsilon\|y_n\|.
\end{align*}
故 $y\in B_n^\varepsilon$。中心为零时对应的球为空，并不妨碍覆盖非零点。
\item 对应\subcref{exercise:ch05-vector-integration}{item:ch05-vector-simple-approximation}：只需用 $j\ge2$。各 $A_{nj}$ 是 Borel 集，彼此不交；因 $1/j<1$，它们的并恰为 $\mathcal Y\setminus\{0\}$，因为前问覆盖非零点，而 $0$ 不在这些球中。故 $E_{nj}$ 可测、彼此不交，其并为 $\{f\ne0\}$。若 $E_{nj}$ 非空，则 $y_n\ne0$，且在其上
\begin{align*}
\|f-y_n\|<j^{-1}\|y_n\|,\qquad
\|f\|>(1-j^{-1})\|y_n\|.
\end{align*}
因此
\begin{align*}
\mu(E_{nj})\le\frac{\|f\|_1}{(1-j^{-1})\|y_n\|}<\infty.
\end{align*}
空的 $E_{nj}$ 当然也有限测度。令 $g_j=y_n$ 于 $E_{nj}$，在 $f=0$ 处令其为零。它为可测的可数值函数，且逐点有
\begin{align*}
\|g_j\|\le\frac j{j-1}\|f\|,\qquad
\|g_j-f\|\le\frac1{j-1}\|f\|.
\end{align*}
所以 $g_j\to f$ 逐点且在 $L^1$ 半范数中收敛。

对固定 $j$，有限截断 $h_{j,N}=\sum_{n=1}^Ny_n\chi_{E_{nj}}$ 属于 $F_{\mathcal Y}$，并且
\begin{align*}
\|g_j-h_{j,N}\|_1=\sum_{n>N}\|y_n\|\mu(E_{nj})\longrightarrow0,
\end{align*}
因为全部非负项之和等于 $\|g_j\|_1<\infty$。选 $N_j$ 使该尾和小于 $2^{-j}$，令 $h_j=h_{j,N_j}$。于是 $\|h_j-f\|_1\le2^{-j}+\|f\|_1/(j-1)\to0$。另一方面，单调收敛定理给
\begin{align*}
\int\sum_{j=2}^{\infty}\|g_j-h_j\|\,\mathrm d\mu
=\sum_{j=2}^{\infty}\|g_j-h_j\|_1<\infty.
\end{align*}
故此误差级数几乎处处有限，特别其项趋零；结合 $g_j\to f$ 逐点得 $h_j\to f$ 几乎处处。重新从 $1$ 编号即为所求序列。
\item 对应\subcref{exercise:ch05-vector-integration}{item:ch05-vector-integral-existence}：先在 $F_{\mathcal Y}$ 上定义。将有限表示细分为两两不交可测集 $E_j$ 上的常值 $y_j$，忽略支集之外的零值部分；这些 $E_j$ 均位于原有限测度支集之并内，因此测度有限。令
\begin{align*}
J\left(\sum_jy_j\chi_{E_j}\right)=\sum_j\mu(E_j)y_j.
\end{align*}
两种表示再取共同有限细分，由有限可加性得到相同的向量和，故良定义；同一论证给线性性。三角不等式给
\begin{align*}
\|Jh\|\le\sum_j\mu(E_j)\|y_j\|=\|h\|_1.
\end{align*}
尤其 $h=k$ 几乎处处时 $Jh=Jk$。

对 $f\in L_{\mathcal Y}^1$，取前问给出的 $h_n\in F_{\mathcal Y}$，满足 $\|h_n-f\|_1\to0$。由 $\|Jh_n-Jh_m\|\le\|h_n-h_m\|_1$，$(Jh_n)$ 在完备的 $\mathcal Y$ 中收敛，定义 $\int f=\lim_nJh_n$。若用另一列 $k_n$，则像之差范数至多 $\|h_n-f\|_1+\|k_n-f\|_1\to0$，所以定义与逼近无关。对 $af+bg$ 使用 $ah_n+bk_n$ 逼近并取极限，得线性性。由 $|\|h_n\|_1-\|f\|_1|\le\|h_n-f\|_1$，取极限还得 $\|\int f\|\le\|f\|_1$。在简单函数上显然有要求的公式。

若另一个线性映射也满足给定公式和界，它在所有简单函数上等于 $J$，且与 $\int$ 一样连续；对任意 $f$ 用 $h_n$ 逼近，差的范数至多 $2\|f-h_n\|_1\to0$，故二者相等。唯一性证明完成。
\item 对应\subcref{exercise:ch05-vector-integration}{item:ch05-vector-dct}：取一个可测零集 $N$，包含收敛和可数多个控制不等式失败的全部点。在 $X\setminus N$ 上取给定极限 $f$，在 $N$ 上定义为零，记这个版本为 $\widetilde f$。同时把各 $f_n$ 在 $N$ 上设零后，得到处处收敛的可测映射列。

此极限可测：对可数稠密集中的每个 $y_j$，距离函数 $\|f_n-y_j\|$ 逐点趋于 $\|\widetilde f-y_j\|$，故后者实值可测。以稠密中心和正有理半径作可数基，每个基球的原像可测；每个开集是这些基球的可数并，所以 $\widetilde f$ Borel 可测。又 $\|\widetilde f\|\le g$ 几乎处处，故属于 $L_{\mathcal Y}^1$。标量控制收敛定理应用于 $\|f_n-\widetilde f\|$，它几乎处处趋零且至多 $2|g|$，得到 $\|f_n-\widetilde f\|_1\to0$。于是
\begin{align*}
\left\|\int f_n-\int\widetilde f\right\|\le\|f_n-\widetilde f\|_1\longrightarrow0.
\end{align*}
若给定的 $f$ 本来可测，则与 $\widetilde f$ 几乎处处相等，积分相同；一般则用这个可测代表解释 $\int f$。
\item 对应\subcref{exercise:ch05-vector-integration}{item:ch05-vector-integral-operator}：$T$ 连续，故 $T\circ f$ 可测；且 $\|T\circ f\|_1\le\|T\|\|f\|_1<\infty$。简单函数 $h=\sum_jy_j\chi_{E_j}$ 满足 $T\circ h=\sum_j(Ty_j)\chi_{E_j}$，所以按定义和线性性有 $\int T\circ h=T(\int h)$。取 $h_n\in F_{\mathcal Y}$ 在 $L^1$ 中逼近 $f$，则 $\|T\circ h_n-T\circ f\|_1\le\|T\|\|h_n-f\|_1\to0$。两边的积分连续性与 $T$ 的连续性使简单函数的恒等式可取极限，得到 $\int T\circ f=T(\int f)$。
\end{enumerate}
\end{solution}
\begin{remark}
\textbf{动机与逐条件说明（重构新增）。}可分性提供可数基和可数稠密中心，分别用于联合可测性及可数分块逼近。半径按中心范数成比例，使非零分块上的 $\|f\|$ 有正常数下界；结合可积性才得到每一块测度有限，而不要求整个测度空间 $\sigma$ 有限。完备性用于简单函数积分的 Cauchy 列在 $\mathcal Y$ 中有极限；唯一性来自范数界和稠密逼近。控制收敛中的单一可积函数同时控制全部项，最后一问中的有界线性算子则保证有限和规则和积分误差都能传过 $T$。
\end{remark}
% END SUPPLEMENT: id=v14-ch05-ex16-solution
\end{document}
